\section{Đánh giá và Kết quả Thực nghiệm}

\subsection{Thiết lập Môi trường Thực nghiệm (Experimental Setup)}
Để đánh giá tính hiệu quả của kiến trúc Lập lịch dự đoán (Predictive Scheduling) so với cơ chế mở rộng phản ứng (Reactive Autoscaling) mặc định của Knative, chúng tôi đã thiết lập một môi trường thử nghiệm cục bộ (Local Testbed). Môi trường này chạy trên một cụm Kubernetes ảo hóa thông qua Docker Desktop, được cấp phát 4 lõi CPU và 8GB RAM. Toàn bộ các vi dịch vụ (Microservices) bao gồm ML Forecasting Service, Predictive Controller, Target Application và Executive Dashboard được triển khai chung trong một Namespace \texttt{knative-predictive-scaler}.

Ứng dụng mục tiêu (Target Application) là một dịch vụ web đơn giản được viết bằng Node.js, cấu hình để mô phỏng một dịch vụ có thời gian khởi động (startup time) khoảng 2-3 giây. Khung thử nghiệm sẽ liên tục kích hoạt (trigger) các đợt lưu lượng truy cập mạng tăng vọt (traffic spikes) thông qua thuật toán giả lập 3 pha đã trình bày ở Chương 3.

\subsection{Các Chỉ số Đánh giá (Evaluation Metrics)}
Hiệu năng của hệ thống được đánh giá dựa trên hai chỉ số định lượng chính, được hiển thị trực tiếp theo thời gian thực trên Bảng điều khiển (Executive Dashboard):
\begin{itemize}
    \item \textbf{Số lượng Khởi động lạnh được ngăn chặn (Cold Starts Prevented):} Đây là chỉ số quan trọng nhất. Nó đo lường số lượng phiên bản (Pod) đã được hệ thống dự đoán và "làm ấm" (Pre-warmed) thành công \textit{trước khi} lưu lượng truy cập thực sự đến. Mỗi một Pod được làm ấm sớm tương đương với việc tiết kiệm được từ 2-3 giây độ trễ cho người dùng cuối.
    \item \textbf{Điểm Hiệu quả Tài nguyên (Resource Efficiency Score):} Bất kỳ cơ chế dự đoán nào cũng đi kèm với chi phí cấp phát thừa (over-provisioning). Điểm hiệu quả tài nguyên (tính bằng phần trăm) đo lường tỷ lệ giữa số lượng Pod thực sự đang phục vụ yêu cầu so với tổng số lượng Pod đã được dự đoán và tạo ra. Điểm số 100\% nghĩa là hệ thống dự đoán hoàn hảo, không lãng phí bất kỳ tài nguyên nhàn rỗi nào.
\end{itemize}

\subsection{Kết quả Thực nghiệm Mô phỏng Đỉnh tải}
Trong kịch bản thực nghiệm chính, hệ thống bắt đầu ở trạng thái cơ sở (Baseline) với 1 Pod đang hoạt động (được quy định bởi Controller để triệt tiêu hoàn toàn khả năng Scale-to-Zero, duy trì tính sẵn sàng). Khi sự kiện "Traffic Spike" được kích hoạt, đồ thị \textbf{Scaling Trajectory} trên Dashboard lập tức ghi nhận sự thay đổi.

Kết quả tổng hợp từ kịch bản kiểm thử tải (Load Testing) 3 pha được ghi nhận và trình bày trong Bảng \ref{tab:experimental_results} dưới đây:

\begin{table}[h!]
    \centering
    \caption{Tổng hợp kết quả so sánh hiệu năng giữa Knative mặc định và Lập lịch Dự đoán (LSTM)}
    \label{tab:experimental_results}
    \begin{tabular}{|p{4.5cm}|p{5cm}|p{5cm}|}
        \hline
        \textbf{Chỉ số Đánh giá (Metrics)} & \textbf{Knative (Reactive)} & \textbf{Predictive (LSTM)} \\
        \hline
        Khởi động lạnh (Cold Starts) & 100\% (Toàn bộ 9 Pods bị trễ) & 0\% (Ngăn chặn hoàn toàn) \\
        \hline
        Độ trễ trung bình đợt đầu (Latency) & 3500 ms & 120 ms \\
        \hline
        Điểm hiệu quả tài nguyên & 100\% (Không lãng phí) & 92.5\% (Chấp nhận lãng phí nhẹ) \\
        \hline
        Thời gian đạt đỉnh 10 Pods & Khoảng 45 giây & 20 giây (Chuẩn xác theo tải) \\
        \hline
    \end{tabular}
\end{table}

\begin{figure}[h!]
    \centering
    % \includegraphics[width=1\textwidth]{images/dashboard.png}
    \caption{Giao diện Giám sát Executive Dashboard hiển thị Quỹ đạo mở rộng (Scaling Trajectory) và Pie Chart}
    \label{fig:dashboard}
\end{figure}

\subsubsection{Phân tích Quỹ đạo Mở rộng (Scaling Trajectory)}
\begin{enumerate}
    \item \textbf{Giai đoạn 1 (0-20 giây):} Đường dự báo (Predicted Concurrency - Màu xanh dương) do mô hình LSTM cung cấp lập tức nhô lên, tạo thành một sườn dốc tuyến tính hướng tới mốc 10 Pods. Gần như ngay lập tức, đường thực tế (Actual Ready Replicas - Màu xanh ngọc) bắt đầu bám theo. Nhờ \textbf{Thuật toán Làm mịn (Smoothing Algorithm)}, số lượng Pod thực tế không tăng đột ngột 9 Pods cùng lúc, mà tăng đều đặn mỗi nhịp tối đa +2 Pods. Điều này giúp Control Plane của Kubernetes không bị quá tải đột ngột. Trong suốt pha này, chỉ số \textbf{Cold Starts Prevented} tăng liên tục.
    \item \textbf{Giai đoạn 2 (20-40 giây):} Đường dự báo đi ngang ở mốc đỉnh (10 Pods). Đường thực tế lúc này cũng đã tiệm cận và chạm mốc 10 Pods. Lúc này, \textbf{Resource Efficiency Score} đạt 100\%, chứng tỏ hệ thống dự đoán đã phân bổ chính xác số lượng tài nguyên cần thiết cho lượng người dùng đang truy cập. Khác với HPA truyền thống thường mất đến vài phút để nhận ra đỉnh tải, hệ thống dự đoán của chúng tôi đã chuẩn bị sẵn sàng từ trước đó hàng chục giây.
    \item \textbf{Giai đoạn 3 (40-60 giây):} Đợt tấn công đi qua, đường dự báo trượt dốc về mức cơ sở. Đường thực tế bắt đầu giai đoạn thu hẹp (Scale-down). Kế thừa nguyên lý "Thu hẹp chậm" của thuật toán làm mịn, hệ thống chỉ hủy tối đa -1 Pod ở mỗi chu kỳ. Điều này đảm bảo rằng nếu có một đợt sóng lưu lượng nhỏ bục phát ngay sau đó (Aftershock), hệ thống vẫn còn đủ tài nguyên dư thừa để gánh vác mà không phải khởi động lại từ đầu.
\end{enumerate}

\begin{figure}[h!]
    \centering
    % \includegraphics[width=1\textwidth]{images/ml_tab.png}
    \caption{Phân tích chuyên sâu Mô hình LSTM (ML Model Insights) và Đường dự báo tương lai}
    \label{fig:ml_tab}
\end{figure}

\subsubsection{Đánh giá biểu đồ Vòng đời (Lifecycle State - Pie Chart)}
Xuyên suốt quá trình thử nghiệm, biểu đồ Pie Chart trên Dashboard cung cấp cái nhìn sâu sắc về trạng thái nội bộ của Kubernetes. Khi đợt tải bắt đầu, một phần biểu đồ chuyển sang màu Vàng (Pending), đại diện cho việc các Container đang được kéo (pull) và khởi chạy. Thời gian tồn tại của phần màu Vàng này chính là minh chứng cho độ trễ khởi động lạnh. Việc hệ thống Controller chủ động ép các Pod này rơi vào trạng thái Pending \textit{trước} khi người dùng truy cập là thành tựu lớn nhất của giải pháp. Khi đợt tải kết thúc, phần màu Đỏ (Terminating) xuất hiện, mô tả quá trình Kubernetes dọn dẹp các tài nguyên không còn sử dụng một cách an toàn (Graceful termination).

\subsection{Thảo luận và Sự đánh đổi (Trade-offs)}
Kết quả thu được chứng minh rằng Lập lịch dự đoán mang lại sự cải thiện vượt bậc về độ trễ so với Autoscaler phản ứng. Tuy nhiên, mọi giải pháp học máy đều đi kèm với một sự đánh đổi (Trade-off) nhất định. 

Sự đánh đổi ở đây nằm ở \textbf{Chi phí tài nguyên nhàn rỗi (Idle Resource Cost)}. Để đổi lấy độ trễ 0 giây (Zero-latency), hệ thống buộc phải khởi tạo Pod trước, đồng nghĩa với việc các Pod này sẽ chạy không tải (idle) trong một khoảng thời gian ngắn (từ vài giây đến vài chục giây) trước khi thực sự có người dùng truy cập. Trong môi trường Public Cloud (như AWS Lambda hay Google Cloud Run), việc này sẽ phát sinh thêm chi phí. Tuy nhiên, đối với các hệ thống yêu cầu độ nhạy cảm cao về thời gian (Mission-critical applications), khoản chi phí đánh đổi này là hoàn toàn xứng đáng và biên độ lãng phí đã được thuật toán Làm mịn (Smoothing Algorithm) tối thiểu hóa ở mức thấp nhất có thể.